
Soit \(\underline{E}\) un \(\underline{A}\)-Module localement libre de rang fini \(n\). Il existe un torseur \(\underline{P}\) sous \(\underline{\GL}(n,\underline{A})\) tel que \(\underline{E}=\underline{P}\times_{\underline{\GL}(n,\underline{A})} (\underline{A})^n\), où \((\underline{A})^n\) est considéré comme \(\underline{\GL}(n,\underline{A})\)-Module par la représentation standard. On peut alors définir un \(\lambda\)-homomorphisme de \(R_{\mathbb Z}(\underline{\underline{\operatorname{Gl}}}(n))\) dans \(K\), en associant à toute représentation \(\rho\) de \(\underline{\underline{\operatorname{Gl}}}(n)\) dans un \(\mathbb Z\)-module libre de rang \(p\) la classe du \(\underline{A}\)-Module localement libre
\(\underline{P}\times_{\underline{\GL}(n,\underline{A})} (\underline{A})^p\), où \((\underline{A})^p\) est considéré comme \(\underline{\GL}(n,\underline{A})\)-module par la représentation \(\rho\). Dans cet homomorphisme, la classe de \(\underline{E}\) est donc l'image de la représentation standard. Pour vérifier une relation polynomiale entre les \(\lambda^i(E)\), il suffit donc que cette relation soit vérifiée par la représentation standard dans \(R_{\mathbb Z}(\underline{\underline{\operatorname{Gl}}}(n))\).

On peut faire un raisonnement analogue avec deux \(\underline{A}\)-Modules localement libres \(\underline{E},\underline{F}\) de rangs \(n,p\) ; on peut alors trouver un torseur \(\underline{P}\) sous \(\underline{\GL}(n,\underline{A})\times\underline{\GL}(p,\underline{A})\) tel que \(\underline{E}=\underline{P}\times_{\underline{\GL}(n,\underline{A})\times\underline{\GL}(p,\underline{A})} (\underline{A})^n\), où \((\underline{A})^n\) est considéré comme \(\underline{\GL}(n,\underline{A})\times\underline{\GL}(p,\underline{A})\)-Module en faisant opérer \(\underline{\GL}(n,\underline{A})\) par la représentation standard et \(\underline{\GL}(p,\underline{A})\) par la représentation triviale, et réciproquement pour \(\underline{F}\). On définit ensuite un \(\lambda\)-homomorphisme de \(R_{\mathbb Z}(\underline{\underline{\operatorname{Gl}}}(n)\times\underline{\underline{\operatorname{Gl}}}(p))\) dans \(K\) de la même façon que plus haut, et on se ramène ainsi à vérifier les relations entre les \(\lambda^i(E)\), \(\lambda^j(F)\) pour les éléments correspondants de \(R_{\mathbb Z}(\underline{\underline{\operatorname{Gl}}}(n)\times\underline{\underline{\operatorname{Gl}}}(p))\).

On est donc ramené à montrer que pour tout schéma en groupes \(G\) sur \(\Spec\mathbb Z\), produit fini de schémas en groupes linéaires, le pré-\(\lambda\)-anneau \(R_{\mathbb Z}(G)\) est un \(\lambda\)-anneau. Pour ce faire, on peut référer à [5], où \(R_{\mathbb Z}(G)\) est calculé explicitement. Si l'on ne veut pas utiliser ce résultat, on peut raisonner de façon analogue à 3.2, en utilisant 1.14.

